\section{Dáta a ich spracovanie}
\label{sec:data}

\subsection{Dataset}

Pre klasifikáciu bol použitý výber z databázy jedál \textbf{Food-101}
(zdroj Kaggle, \texttt{food41}). Pôvodná databáza obsahuje 101 tried, podľa
zadania (,,použiť minimálne 10 tried, ktoré sú najviac zastúpené``) bolo
vybraných \textbf{15 tried}: \emph{beef\_tartare, chicken\_wings, cup\_cakes,
donuts, dumplings, french\_fries, greek\_salad, hot\_dog, lasagna,
lobster\_bisque, pancakes, panna\_cotta, pizza, steak, waffles}. Každá trieda
obsahuje presne 1000 obrázkov, spolu \textbf{15\,000 obrázkov}, dataset je
teda dokonale vyvážený.

\subsection{Rozdelenie na trénovaciu, validačnú a testovaciu množinu}

Dataset bol náhodne rozdelený v pomere \textbf{80\,\%~/~10\,\%~/~10\,\%}
(stratifikovane podľa triedy, aby malo každé jedlo rovnaké zastúpenie vo
všetkých troch množinách):

\begin{table}[H]
    \centering
    \caption{Veľkosť jednotlivých množín}
    \label{tab:data}
    \begin{tabular}{lcc}
        \toprule
        Množina & Obrázkov spolu & Obrázkov na triedu \\
        \midrule
        Trénovacia  & 12\,000 & 800 \\
        Validačná   & 1\,500  & 100 \\
        Testovacia  & 1\,500  & 100 \\
        \bottomrule
    \end{tabular}
\end{table}

Rozdelenie je uložené v \texttt{data\_split/manifest.csv} s pevne daným
náhodným semienkom, takže je reprodukovateľné. Všetky modely a konfigurácie
používajú to isté rozdelenie a testovacie dáta neboli pri trénovaní ani pri
výbere najlepšej epochy použité (najlepší checkpoint sa vyberá podľa
validačnej accuracy).

\subsection{Predspracovanie a augmentácia}

Všetky obrázky sú pred vstupom do siete zmenšené a orezané na $224\times224$
px a normalizované strednou hodnotou a smerodajnou odchýlkou datasetu
ImageNet (keďže backbone je naň predtrénovaný). Pre validačnú a testovaciu
množinu sa používa iba táto deterministická transformácia.

Pre trénovaciu množinu je k dispozícii augmentovaná verzia
(\texttt{RandomResizedCrop} so škálou 0,7--1,0, horizontálne preklopenie,
náhodná rotácia $\pm15^\circ$, farebný jitter, náhodný autokontrast
a \texttt{RandomErasing}), ktorá sa používa v základnej konfigurácii oboch
modelov a je vypnutá pri experimente bez augmentácie
(kapitola~\ref{sec:augmentacia}).

\subsection{Parametre trénovania}
\label{sec:params}

Podľa zadania sa pri všetkých modeloch aj konfiguráciách používajú rovnaké
parametre trénovania, takže rozdiely vo výsledkoch možno pripísať
architektúre, resp. skúmanému faktoru:

\begin{table}[H]
    \centering
    \caption{Parametre trénovania (spoločné pre všetky modely M1--M5)}
    \label{tab:params}
    \begin{tabular}{ll}
        \toprule
        Parameter & Hodnota \\
        \midrule
        Optimalizátor & AdamW \\
        Learning rate & $3 \times 10^{-4}$ \\
        Weight decay & $1 \times 10^{-4}$ \\
        Veľkosť dávky & 128 \\
        Počet epoch & 30 \\
        Strata & CrossEntropyLoss \\
        Presnosť výpočtu & zmiešaná (AMP, float16) \\
        Výber checkpointu & epocha s najvyššou validačnou accuracy \\
        Počet nezávislých behov & 3 (semienka 43, 44, 45) \\
        \bottomrule
    \end{tabular}
\end{table}

AdamW s learning rate $3 \times 10^{-4}$ je bežná a robustná voľba pre
transferové učenie -- dostatočne malý, aby pri rozmrazení vrstiev
nepoškodil predtrénované váhy, a zároveň dostatočne veľký, aby sa nová
výstupná vrstva naučila za 30 epoch. Veľkosť dávky 128 bola zvolená ako
maximum, ktoré sa pri postupnom rozmrazovaní M2 (najnáročnejšia
konfigurácia, $\sim$5,8~GiB) zmestí do pamäte použitej GPU (8~GiB).

Modely M3--M5 by po úplnom rozmrazení pri dávke 128 potrebovali 7,7~GiB
(ConvNeXt-T), 8,7~GiB (Swin-T) a 10,6~GiB (CoAtNet-0), čo je na hranici
alebo nad kapacitou GPU. Aby zostala veľkosť dávky rovnaká pre všetky
modely, pri ich postupnom rozmrazovaní bol zapnutý \textbf{gradient
checkpointing} -- aktivácie vnútri stupňov sa počas dopredného prechodu
neukladajú, ale pri spätnom prechode sa prepočítajú znova. Pamäť klesne na
2,3--3,3~GiB, výsledok výpočtu sa nemení, trénovanie je však pomalšie.
Jediný parameter, ktorý sa medzi modelmi líši, je normalizácia vstupu: musí
zodpovedať štatistikám, s ktorými bol model predtrénovaný (ImageNet pre
M1--M4, stredná hodnota aj smerodajná odchýlka 0,5 pre CoAtNet v M5).

Metriky Precision, Recall a F1 sú počítané ako \emph{makro} priemer cez
triedy. Keďže testovacia množina má pre každú triedu presne 100 obrázkov,
makro Recall je číselne zhodný s Accuracy. Smerodajná odchýlka je všade
počítaná z troch behov (populačná, $\sigma$).

\section{Model M1 -- CNN ResNet50}
\label{sec:m1}

Ako architektúra CNN bol zvolený \textbf{ResNet50}, predtrénovaný na
ImageNet (váhy \texttt{IMAGENET1K\_V2}) -- osvedčená architektúra s
reziduálnymi spojeniami, ktorá je dostatočne výkonná, a pritom sa (na
rozdiel od väčších modelov) dá pri dávke 128 trénovať na bežnej GPU.
Výstupná plne prepojená vrstva bola nahradená novou lineárnou vrstvou
s 15 výstupmi. V základnej konfigurácii je celý backbone \textbf{zamrznutý}
-- trénuje sa iba nová výstupná vrstva (30\,735 parametrov z 23,5 milióna).
Táto konfigurácia je referenčná pre porovnanie s modelom M2
aj s postupným rozmrazovaním.

\subsection{Priebeh trénovania a výsledky na testovacích dátach}

Obr.~\ref{fig:m1-curves} zobrazuje priebeh trénovania najlepšieho z troch
behov. Krivky rastú plynulo a trénovacia a validačná krivka sa takmer
prekrývajú -- model sa nepreučuje. Keďže sa trénuje iba lineárna vrstva nad
zamrznutými príznakmi, model má vzhľadom na 12\,000 trénovacích obrázkov
malú kapacitu a jeho úspešnosť je obmedzená skôr kvalitou (všeobecných,
ImageNet) príznakov než množstvom dát.

\obrazok[0.9\textwidth]{img/m1_curves.png}{M1 -- priebeh trénovania
najlepšieho behu (Loss a Accuracy)}{fig:m1-curves}

\begin{table}[H]
    \centering
    \caption{M1 -- testovacie metriky (priemer $\pm$ smerodajná odchýlka z 3 behov)}
    \label{tab:m1-metrics}
    \begin{tabular}{lc}
        \toprule
        Metrika & Hodnota \\
        \midrule
        Accuracy  & 86,13 $\pm$ 0,16~\% \\
        Precision & 86,43 $\pm$ 0,15~\% \\
        Recall    & 86,13 $\pm$ 0,16~\% \\
        F1        & 86,19 $\pm$ 0,16~\% \\
        \bottomrule
    \end{tabular}
\end{table}

Nízka smerodajná odchýlka (0,16~p.\,b.) ukazuje, že výsledok takmer
nezávisí od náhodnej inicializácie -- pri zamrznutom backbone sa medzi
behmi líši iba inicializácia jednej lineárnej vrstvy a poradie dávok.

\obrazok[0.85\textwidth]{img/m1_confusion.png}{M1 -- kontingenčná tabuľka
najlepšieho behu (testovacie dáta)}{fig:m1-confusion}

Súčtom kontingenčných tabuliek všetkých troch behov (300 testovacích
obrázkov na triedu) vychádzajú ako najčastejšie zámeny \emph{lasagna}
$\rightarrow$ \emph{pizza} (19-krát), \emph{chicken\_wings} $\rightarrow$
\emph{steak} (17), \emph{beef\_tartare} $\rightarrow$ \emph{steak} (16)
a \emph{pizza} $\rightarrow$ \emph{lasagna} (15). Ide o vizuálne príbuzné
jedlá -- zapečený syr a paradajková omáčka, resp. hnedé opečené mäso.
Najhoršie rozpoznávané triedy sú \emph{pancakes} (recall 80,3~\%),
\emph{pizza} (80,7~\%) a \emph{beef\_tartare} (82,3~\%), najlepšie
\emph{dumplings} (91,7~\%), \emph{greek\_salad} (91,0~\%)
a \emph{lobster\_bisque} (89,7~\%).

\subsection{Vizualizácia rozhodovania -- GradCAM}

Pre náhodne vybrané testovacie vzorky (4 správne a 4 nesprávne
klasifikované) bola vykonaná vizualizácia rozhodovania siete metódou
\textbf{GradCAM}, aplikovanou na výstup posledného konvolučného bloku
(\texttt{layer4}). Mapa zvýrazňuje oblasti obrazu, ktoré najviac prispeli
k predikovanej triede.

\obrazok[\textwidth]{img/m1_gradcam.png}{M1 -- GradCAM vizualizácia na
správne (hore) a nesprávne (dole) klasifikovaných testovacích
vzorkách}{fig:m1-gradcam}

Pri správne klasifikovaných vzorkách sa aktivácia sústreďuje na samotné
jedlo a jeho typickú textúru. Pri chybách sieť často reaguje na časť obrazu,
ktorá je vizuálne zameniteľná s inou triedou, alebo na okolie (tanier,
prílohu), čo zodpovedá zámenám v kontingenčnej tabuľke.

\section{Model M2 -- ResNet50 + CBAM}
\label{sec:m2}

\subsection{Architektúra}

Model M2 je podľa zadania model M1 doplnený o attention mechanizmus
\textbf{CBAM} (Convolutional Block Attention Module, Woo a kol., 2018).
CBAM postupne aplikuje dva typy pozornosti na mapu príznakov
$F \in \mathbb{R}^{C \times H \times W}$:

\begin{itemize}
    \item \textbf{kanálová pozornosť} -- príznaky sa zredukujú priemerným
          aj maximálnym poolingom cez priestorové rozmery, prejdú zdieľaným
          MLP (redukcia $C \rightarrow C/16 \rightarrow C$) a súčet oboch
          výstupov po sigmoide váhuje jednotlivé kanály -- sieť sa učí,
          \emph{ktoré} príznaky sú dôležité,
    \item \textbf{priestorová pozornosť} -- priemer a maximum cez kanály
          sa spoja do dvojkanálovej mapy, konvolúcia $7\times7$ so sigmoidom
          z nej vytvorí váhu pre každú pozíciu -- sieť sa učí, \emph{kde}
          v obraze sú dôležité oblasti.
\end{itemize}

\noindent Blok CBAM je vložený za každý zo štyroch reziduálnych stupňov
ResNet50 (\texttt{layer1}--\texttt{layer4}, 256, 512, 1024 a 2048 kanálov).
Predtrénovaný backbone aj výstupná vrstva sú rovnaké ako pri M1.

Pri zamrznutom backbone sa v M2 trénujú bloky CBAM aj výstupná vrstva --
\textbf{731\,527 trénovateľných parametrov} oproti 30\,735 pri M1 (celkovo
24,2 mil. oproti 23,5 mil.). Keďže CBAM je vložený \emph{do} backbonu, gradient
musí prejsť spätne aj cez zamrznuté vrstvy \texttt{layer2}--\texttt{layer4}
(ich váhy sa nemenia, ale treba vypočítať gradient vzhľadom na aktivácie).
Preto je jedna epocha M2 približne dvakrát pomalšia než pri M1 a model
potrebuje viac pamäte GPU.

\subsection{Priebeh trénovania a výsledky na testovacích dátach}

\obrazok[0.9\textwidth]{img/m2_curves.png}{M2 -- priebeh trénovania
najlepšieho behu (Loss a Accuracy)}{fig:m2-curves}

Na rozdiel od M1 začína M2 pomaly -- po prvej epoche má validačnú accuracy
len $\sim$15--50~\%. Náhodne inicializované bloky CBAM totiž na začiatku
náhodne preváhujú predtrénované príznaky, a model sa najskôr musí naučiť
pozornosť, ktorá ich nepoškodzuje. Približne od 4. epochy sa M2 vyrovná M1
a potom ho predbieha. Na rozdiel od M1 sa pri M2 postupne otvára rozdiel
medzi trénovacou (93,6~\%) a validačnou (88,1~\%) accuracy -- väčší počet
trénovateľných parametrov zvyšuje kapacitu, ale aj náchylnosť na preučenie.

\begin{table}[H]
    \centering
    \caption{M2 -- testovacie metriky (priemer $\pm$ smerodajná odchýlka z 3 behov)}
    \label{tab:m2-metrics}
    \begin{tabular}{lc}
        \toprule
        Metrika & Hodnota \\
        \midrule
        Accuracy  & 89,07 $\pm$ 0,67~\% \\
        Precision & 89,33 $\pm$ 0,64~\% \\
        Recall    & 89,07 $\pm$ 0,67~\% \\
        F1        & 89,09 $\pm$ 0,67~\% \\
        \bottomrule
    \end{tabular}
\end{table}

Jednotlivé behy dosiahli 88,73~\%, 90,00~\% a 88,47~\%. Rozptyl je približne
štyrikrát väčší ako pri M1 -- s viac ako 700 tisíc náhodne
inicializovanými parametrami CBAM závisí výsledok viac od inicializácie.

\obrazok[0.85\textwidth]{img/m2_confusion.png}{M2 -- kontingenčná tabuľka
najlepšieho behu (testovacie dáta)}{fig:m2-confusion}

V súčte troch behov sú najčastejšie zámeny \emph{lasagna} $\rightarrow$
\emph{pizza} (20-krát, rovnako ako pri M1), \emph{panna\_cotta}
$\rightarrow$ \emph{lobster\_bisque} (17) a \emph{cup\_cakes}
$\rightarrow$ \emph{donuts} (14). Zámeny mäsových jediel za \emph{steak},
typické pre M1, sa výrazne zmenšili. Najhoršie rozpoznávaná trieda je
\emph{steak} (85,0~\%), najlepšie \emph{greek\_salad} (95,7~\%)
a \emph{dumplings} (94,7~\%).

\subsection{Vizualizácia rozhodovania -- GradCAM}

Pri M2 je GradCAM aplikovaný na výstup posledného bloku CBAM
(\texttt{cbam4}), teda na príznaky \texttt{layer4} po zohľadnení pozornosti --
práve tie vstupujú do klasifikátora.

\obrazok[\textwidth]{img/m2_gradcam.png}{M2 -- GradCAM vizualizácia na
správne (hore) a nesprávne (dole) klasifikovaných testovacích
vzorkách}{fig:m2-gradcam}

Pri správnych predikciách sa aktivácia sústreďuje na charakteristické
časti jedla (mriežka wafle, poleva cupcaku). Chyby sú dobre
interpretovateľné: pri hot dogu klasifikovanom ako \emph{french\_fries}
sa model pozerá na hranolky položené na párku, pri šiškách
klasifikovaných ako \emph{cup\_cakes} na zdobenú polevu, ktorá je typická
pre obe triedy. Chyby teda vznikajú hlavne pri obrázkoch, ktoré naozaj
obsahujú prvky inej triedy.

\section{Vplyv augmentácie trénovacích dát}
\label{sec:augmentacia}

Podľa zadania (,,aspoň pre jeden model ukážte vplyv augmentácie na
úspešnosť klasifikácie``) bola pre oba modely porovnaná základná
konfigurácia s augmentáciou a konfigurácia bez augmentácie. Ostatné
parametre trénovania (Tabuľka~\ref{tab:params}) zostali nezmenené.

\begin{table}[H]
    \centering
    \caption{Vplyv augmentácie -- testovacie metriky (priemer $\pm$ std z 3 behov)}
    \label{tab:aug-table}
    \footnotesize
    \begin{tabular}{lcccc}
        \toprule
         & \multicolumn{2}{c}{M1} & \multicolumn{2}{c}{M2} \\
        \cmidrule(lr){2-3} \cmidrule(lr){4-5}
         & s augm. & bez augm. & s augm. & bez augm. \\
        \midrule
        Accuracy  & \textbf{86,13 $\pm$ 0,16} & 84,31 $\pm$ 0,06 & \textbf{89,07 $\pm$ 0,67} & 85,53 $\pm$ 0,52 \\
        Precision & \textbf{86,43 $\pm$ 0,15} & 84,67 $\pm$ 0,08 & \textbf{89,33 $\pm$ 0,64} & 86,01 $\pm$ 0,51 \\
        Recall    & \textbf{86,13 $\pm$ 0,16} & 84,31 $\pm$ 0,06 & \textbf{89,07 $\pm$ 0,67} & 85,53 $\pm$ 0,52 \\
        F1        & \textbf{86,19 $\pm$ 0,16} & 84,38 $\pm$ 0,07 & \textbf{89,09 $\pm$ 0,67} & 85,57 $\pm$ 0,48 \\
        \bottomrule
    \end{tabular}
\end{table}

Augmentácia zlepšila testovaciu accuracy pri oboch modeloch -- pri M1
o \textbf{1,8~p.\,b.}, pri M2 o \textbf{3,5~p.\,b.} V oboch prípadoch je rozdiel
niekoľkonásobne väčší než smerodajná odchýlka, ide teda o systematický efekt.

\obrazok[0.75\textwidth]{img/m1_compare_augmentation_metrics.png}{M1 -- vplyv
augmentácie na testovacie metriky}{fig:m1-aug-metrics}

\obrazok[0.9\textwidth]{img/m1_compare_augmentation_curves.png}{M1 -- vplyv
augmentácie na priebeh validačného Loss a Accuracy (najlepší beh)}{fig:m1-aug-curves}

\obrazok[0.75\textwidth]{img/m2_compare_augmentation_metrics.png}{M2 -- vplyv
augmentácie na testovacie metriky}{fig:m2-aug-metrics}

\obrazok[0.9\textwidth]{img/m2_compare_augmentation_curves.png}{M2 -- vplyv
augmentácie na priebeh validačného Loss a Accuracy (najlepší beh)}{fig:m2-aug-curves}

Pri M1 sú validačné krivky s augmentáciou aj bez nej takmer totožné
(Obr.~\ref{fig:m1-aug-curves}) a rozdiel sa prejaví až na testovacích dátach.
M1 má iba 30 tisíc trénovateľných parametrov a nepreučuje sa ani bez
augmentácie, preto augmentácia nepôsobí ako regularizácia proti preučeniu,
ale robí naučený klasifikátor odolnejším voči variáciám vstupu (orezanie,
uhol, osvetlenie).

Pri M2 je efekt dvojnásobný. M2 má vďaka blokom CBAM oveľa väčšiu
kapacitu a bez augmentácie sa rýchlejšie prispôsobí konkrétnym trénovacím
obrázkom -- augmentácia tu teda plní aj úlohu klasickej regularizácie.
Z toho vyplýva, že väčší model (alebo model s viac trénovateľnými
parametrami) z augmentácie profituje viac.

\obrazok[0.8\textwidth]{img/m1_noaug_confusion.png}{M1 bez augmentácie --
kontingenčná tabuľka najlepšieho behu (testovacie dáta)}{fig:m1-noaug-confusion}

\obrazok[0.8\textwidth]{img/m2_noaug_confusion.png}{M2 bez augmentácie --
kontingenčná tabuľka najlepšieho behu (testovacie dáta)}{fig:m2-noaug-confusion}

\section{Vplyv postupného rozmrazovania vrstiev}
\label{sec:unfreeze}

Podľa zadania (,,aspoň pre jeden model realizujte postupné rozmrazovanie
vrstiev predtrénovanej siete``) bola pre oba modely porovnaná základná
konfigurácia so \textbf{zamrznutým backbonom} a konfigurácia
s \textbf{postupným rozmrazovaním}.

Rozmrazovanie prebieha v poradí \texttt{layer4} $\rightarrow$
\texttt{layer3} $\rightarrow$ \texttt{layer2} $\rightarrow$ \texttt{layer1},
teda od najvyšších, najviac úlohovo-špecifických vrstiev k najnižším,
najvšeobecnejším (hrany, textúry). Prvých 6 epoch sa trénuje iba nová
výstupná vrstva (a pri M2 bloky CBAM), aby náhodne inicializované vrstvy
neposlali do backbonu veľké, nezmyselné gradienty. Potom sa každých
6 epoch rozmrazí jeden stupeň: \texttt{layer4} v 7. epoche (+15,0 mil.
parametrov), \texttt{layer3} v 13. (+7,1 mil.), \texttt{layer2} v 19.
(+1,2 mil.) a \texttt{layer1} v 25. epoche (+0,2 mil.). Vstupná konvolúcia
(\emph{stem}) zostáva zamrznutá. Novo rozmrazené parametre sa pridajú do
optimalizátora s rovnakým learning rate.

\begin{table}[H]
    \centering
    \caption{Vplyv postupného rozmrazovania -- testovacie metriky (priemer $\pm$ std z 3 behov)}
    \label{tab:unfreeze-table}
    \footnotesize
    \begin{tabular}{lcccc}
        \toprule
         & \multicolumn{2}{c}{M1} & \multicolumn{2}{c}{M2} \\
        \cmidrule(lr){2-3} \cmidrule(lr){4-5}
         & zamrznutý & rozmrazovanie & zamrznutý & rozmrazovanie \\
        \midrule
        Accuracy  & 86,13 $\pm$ 0,16 & \textbf{92,60 $\pm$ 0,19} & 89,07 $\pm$ 0,67 & \textbf{91,91 $\pm$ 0,85} \\
        Precision & 86,43 $\pm$ 0,15 & \textbf{92,75 $\pm$ 0,15} & 89,33 $\pm$ 0,64 & \textbf{92,05 $\pm$ 0,79} \\
        Recall    & 86,13 $\pm$ 0,16 & \textbf{92,60 $\pm$ 0,19} & 89,07 $\pm$ 0,67 & \textbf{91,91 $\pm$ 0,85} \\
        F1        & 86,19 $\pm$ 0,16 & \textbf{92,61 $\pm$ 0,17} & 89,09 $\pm$ 0,67 & \textbf{91,93 $\pm$ 0,84} \\
        \bottomrule
    \end{tabular}
\end{table}

Postupné rozmrazovanie je v tomto projekte najúčinnejšie opatrenie --
zvýšilo testovaciu accuracy M1 o \textbf{6,5~p.\,b.} a M2
o \textbf{2,8~p.\,b.}

\obrazok[0.75\textwidth]{img/m1_compare_gradual_unfreeze_metrics.png}{M1 -- vplyv
postupného rozmrazovania na testovacie metriky}{fig:m1-unfreeze-metrics}

\obrazok[0.9\textwidth]{img/m1_compare_gradual_unfreeze_curves.png}{M1 -- vplyv
postupného rozmrazovania na priebeh validačného Loss a Accuracy (najlepší
beh)}{fig:m1-unfreeze-curves}

\obrazok[0.75\textwidth]{img/m2_compare_gradual_unfreeze_metrics.png}{M2 -- vplyv
postupného rozmrazovania na testovacie metriky}{fig:m2-unfreeze-metrics}

\obrazok[0.9\textwidth]{img/m2_compare_gradual_unfreeze_curves.png}{M2 -- vplyv
postupného rozmrazovania na priebeh validačného Loss a Accuracy (najlepší
beh)}{fig:m2-unfreeze-curves}

Priebeh validačnej accuracy (Obr.~\ref{fig:m1-unfreeze-curves}
a~\ref{fig:m2-unfreeze-curves}) ukazuje pri oboch modeloch skok
v 7. epoche, keď sa rozmrazí \texttt{layer4}. Pri M1 klesne validačný
loss z $\sim$0,9 na $\sim$0,35 a accuracy stúpne z $\sim$80~\% na
$\sim$90~\%. Ďalšie stupne (\texttt{layer3}--\texttt{layer1}) prinášajú už
len malé zlepšenie. Rozhodujúce je teda doučiť najvyšší konvolučný stupeň,
ktorého príznaky sú najviac špecifické pre pôvodnú úlohu (ImageNet)
a najviac profitujú z prispôsobenia na doménu jedál.

Pri M2 je prínos menší. Jedna časť vysvetlenia je, že bloky CBAM už
v zamrznutej konfigurácii čiastočne prispôsobujú príznaky backbonu úlohe
(prevážením kanálov a oblastí), takže rozmrazenie má menej čo doháňať.
Druhá časť je preučenie: po rozmrazení dosahuje M2 trénovaciu accuracy
$\sim$98~\%, zatiaľ čo validačná sa zastaví na $\sim$90~\%. Validačný loss
po 8. epoche už neklesá a iba kolíše medzi $\sim$0,41 a 0,47.

\obrazok[0.8\textwidth]{img/m1_gradual_confusion.png}{M1 s postupným
rozmrazovaním -- kontingenčná tabuľka najlepšieho behu (testovacie
dáta)}{fig:m1-gradual-confusion}

\obrazok[0.8\textwidth]{img/m2_gradual_confusion.png}{M2 s postupným
rozmrazovaním -- kontingenčná tabuľka najlepšieho behu (testovacie
dáta)}{fig:m2-gradual-confusion}

\section{Porovnanie modelov M1 a M2}
\label{sec:m1m2}

Tabuľka~\ref{tab:m1m2} porovnáva oba modely vo všetkých troch
konfiguráciách. Štatistická významnosť rozdielu testovacej accuracy bola
overená \textbf{Welchovým t-testom} (dvojvýberový t-test bez predpokladu
zhodných rozptylov) nad tromi behmi každého modelu. Pri iba troch behoch
má test malú silu, preto hodnotu $p$ treba chápať ako orientačnú.

\begin{table}[H]
    \centering
    \caption{Porovnanie M1 a M2 -- testovacia accuracy (priemer $\pm$ std z 3 behov) a~Welchov t-test}
    \label{tab:m1m2}
    \footnotesize
    \begin{tabular}{lccccc}
        \toprule
        Konfigurácia & M1 & M2 & Rozdiel M2$-$M1 & $p$ & Významný ($\alpha=0{,}05$) \\
        \midrule
        Základná (s augm.)      & 86,13 $\pm$ 0,16 & \textbf{89,07 $\pm$ 0,67} & +2,93 & 0,020 & áno \\
        Bez augmentácie         & 84,31 $\pm$ 0,06 & \textbf{85,53 $\pm$ 0,52} & +1,22 & 0,078 & nie \\
        Postupné rozmrazovanie  & \textbf{92,60 $\pm$ 0,19} & 91,91 $\pm$ 0,85 & $-$0,69 & 0,369 & nie \\
        \bottomrule
    \end{tabular}
\end{table}

\obrazok[0.85\textwidth]{img/m1m2_accuracy.png}{M1 vs M2 -- testovacia
accuracy vo všetkých troch konfiguráciách (priemer 3 behov $\pm$
std)}{fig:m1m2-acc}

\obrazok[0.75\textwidth]{img/m1m2_metrics_baseline.png}{M1 vs M2 (základná
konfigurácia) -- Accuracy, Precision, Recall a F1}{fig:m1m2-metrics}

\obrazok[0.95\textwidth]{img/m1m2_curves_baseline.png}{M1 vs M2 (základná
konfigurácia) -- priebeh trénovania najlepšieho behu (prerušovane
trénovacie, plne validačné dáta)}{fig:m1m2-curves}

\obrazok[\textwidth]{img/m1m2_class_recall.png}{M1 vs M2 (základná
konfigurácia) -- recall jednotlivých tried, súčet troch behov (testovacie
dáta)}{fig:m1m2-recall}

\subsection{Diskusia}

\paragraph{Zamrznutý backbone.} Pri zamrznutom backbone je M2 štatisticky
významne lepší ako M1 (+2,9~p.\,b., $p = 0{,}020$). Zlepšenie je plošné --
M2 má vyšší recall pri 12 z 15 tried, pri jednej rovnaký a pri dvoch o 1 obrázok nižší
(Obr.~\ref{fig:m1m2-recall}). Najviac pomohol pri triedach, kde M1 zlyhával
najviac: \emph{pizza} (80,7 $\rightarrow$ 90,0~\%), \emph{beef\_tartare}
(82,3 $\rightarrow$ 89,3~\%) a \emph{pancakes} (80,3 $\rightarrow$ 86,7~\%).
Pri \emph{french\_fries}, \emph{steak} a \emph{waffles} sú modely prakticky
rovnaké. Pri interpretácii však treba zohľadniť, že M2 nemá len
attention mechanizmus, ale aj 24-krát viac trénovateľných parametrov
(731 tisíc oproti 30 tisícom). M1 v tejto konfigurácii môže meniť iba
lineárnu vrstvu, M2 môže cez CBAM ovplyvniť aj vnútorné príznaky backbonu.
Časť zlepšenia preto pochádza zo samotnej väčšej prispôsobiteľnosti, nie
nutne z mechanizmu pozornosti ako takého.

\paragraph{Bez augmentácie.} Bez augmentácie je M2 stále lepší
(+1,2~p.\,b.), rozdiel však nie je štatisticky významný ($p = 0{,}078$).
Väčšia kapacita M2 sa bez augmentácie čiastočne minie na prispôsobenie sa
trénovacím obrázkom.

\paragraph{Postupné rozmrazovanie.} Keď sa rozmrazí aj backbone, oba modely
majú k dispozícii desiatky miliónov trénovateľných parametrov a výhoda
M2 zanikne -- M1 je dokonca o 0,7~p.\,b. lepší, rozdiel však nie je
významný ($p = 0{,}37$). Pri plnom doladení siete teda CBAM na tomto
datasete neprináša merateľné zlepšenie, zato zvyšuje náchylnosť na
preučenie a rozptyl výsledkov.

\paragraph{Stabilita a výpočtová náročnosť.} M2 má vo všetkých
konfiguráciách 4- až 9-krát väčšiu smerodajnú odchýlku než M1 (0,5--0,85
oproti 0,06--0,19~p.\,b.). Výsledok M2 teda závisí viac od náhodnej
inicializácie blokov CBAM. Trénovanie M2 trvalo približne 2,2-krát dlhšie
(Tabuľka~\ref{tab:summary}), pretože gradient musí prechádzať cez backbone
aj pri zamrznutých váhach. Časy sú orientačné -- počas trénovania M2 bola
GPU čiastočne zdieľaná s inými procesmi, čo je vidieť aj na rozdielnych
časoch jednotlivých behov postupného rozmrazovania (23 až 39 minút).

\section{Výber predtrénovaných modelov M3--M5}
\label{sec:vyber}

Pre modely M3--M5 boli z knižnice \textbf{timm} (PyTorch Image Models)
vybrané architektúry, ktoré spĺňajú dve podmienky. Po prvé, sú
predtrénované \textbf{iba na ImageNet-1k}, teda na rovnakých dátach ako
ResNet50 v M1 a M2 -- väčšie predtrénovacie datasety (napr. ImageNet-22k)
by modelom dali výhodu, ktorá nesúvisí s architektúrou. Po druhé, majú
\textbf{podobnú veľkosť a výpočtovú náročnosť} ako ResNet50 (verzie
,,Tiny`` resp. ,,0``), takže rozdiely vo výsledkoch možno pripísať
architektúre, nie veľkosti modelu.

\begin{table}[H]
    \centering
    \caption{Porovnanie zvolených predtrénovaných modelov (počet parametrov pri 1000 triedach ImageNet)}
    \label{tab:vyber}
    \footnotesize
    \begin{tabular}{lllccc}
        \toprule
        Model & Typ & Predtrénované váhy & Parametre & GMAC & ImageNet top-1 \\
        \midrule
        M1 ResNet50    & CNN & \texttt{IMAGENET1K\_V2} (torchvision) & 25,6 mil. & 4,1 & 80,9~\% \\
        M3 ConvNeXt-T  & optimalizovaná CNN & \texttt{convnext\_tiny.fb\_in1k} & 28,6 mil. & 4,5 & 82,1~\% \\
        M4 Swin-T      & vizuálny transformer & \texttt{swin\_tiny\_patch4\_window7\_224.ms\_in1k} & 28,3 mil. & 4,5 & 81,2~\% \\
        M5 CoAtNet-0   & hybridný model & \texttt{coatnet\_0\_rw\_224.sw\_in1k} & 27,4 mil. & 4,4 & 82,4~\% \\
        \bottomrule
    \end{tabular}
\end{table}

\begin{itemize}
    \item \textbf{M3 -- ConvNeXt-T} (Liu a kol., \emph{A ConvNet for the
          2020s}, 2022) vznikol postupnou modernizáciou ResNet50 podľa
          návrhových princípov vizuálnych transformerov: vstupný
          ,,patchify`` stem $4\times4$, invertované bottleneck bloky,
          hĺbkové (depthwise) konvolúcie $7\times7$ s veľkým receptívnym
          poľom, LayerNorm namiesto BatchNorm a aktivácia GELU. Je to stále
          čistá CNN, ktorá pri rovnakej náročnosti dosahuje úroveň
          transformerov. Verzia Tiny je priamym náprotivkom ResNet50.
    \item \textbf{M4 -- Swin-T} (Liu a kol., \emph{Swin Transformer}, 2021)
          je hierarchický vizuálny transformer. Self-attention počíta iba
          v lokálnych oknách $7\times7$ tokenov, ktoré sa medzi
          susednými blokmi posúvajú (\emph{shifted windows}), takže
          informácia sa šíri cez hranice okien a zložitosť rastie lineárne
          s veľkosťou obrazu. Medzi stupňami sa rozlíšenie zmenšuje
          (56 $\rightarrow$ 28 $\rightarrow$ 14 $\rightarrow$ 7), čím vzniká
          pyramída príznakov podobná CNN. Zo zadaním ponúkaných možností
          (ViT, DeiT, Swin) bol Swin-T zvolený preto, že má rovnakú veľkosť
          ako ConvNeXt-T a CoAtNet-0 (autori ConvNeXt ho používajú ako
          priamu referenciu) a jeho 4 stupne umožňujú rovnaký rozvrh
          postupného rozmrazovania ako pri ResNet50. Pôvodný ViT-B/16
          (86~mil. parametrov, 17,6~GMAC) je približne 4-krát náročnejší a po
          rozmrazení by sa pri dávke 128 nezmestil do 8~GiB GPU, DeiT-S
          (22~mil., 79,9~\%) je čistý ViT s nižšou presnosťou na ImageNet.
    \item \textbf{M5 -- CoAtNet-0} (Dai a kol., \emph{CoAtNet: Marrying
          Convolution and Attention}, 2021) spája konvolúciu a pozornosť
          ,,vertikálne``: prvé dva stupne tvoria konvolučné bloky MBConv
          (hĺbkovo separovateľné konvolúcie s BatchNorm), ktoré dodávajú
          induktívny predpoklad lokality vhodný pre menšie dáta, posledné
          dva stupne sú transformerové bloky s relatívnou pozornosťou
          s globálnym receptívnym poľom. Keďže pôvodné váhy autorov nie sú
          verejne dostupné, bola použitá implementácia a váhy z timm
          (variant \texttt{rw}). Tento model bol predtrénovaný s normalizáciou
          vstupu so strednou hodnotou aj smerodajnou odchýlkou 0,5, ktorá
          bola preto použitá aj pri trénovaní a GradCAM.
\end{itemize}

\subsection{Spoločný postup pre M3--M5}
\label{sec:m345-postup}

Všetky tri modely sú vytvorené rovnako ako M1: predtrénované váhy, nová
lineárna výstupná vrstva s 15 výstupmi (768 vstupov, 11\,535 parametrov)
a v základnej konfigurácii \textbf{zamrznutý backbone}, v ktorom sa trénuje
iba táto nová vrstva. Druhá konfigurácia, \textbf{postupné rozmrazovanie},
používa rovnaký rozvrh ako M1 a M2 (kapitola~\ref{sec:unfreeze}): prvých
6 epoch sa trénuje iba výstupná vrstva, potom sa v 7., 13., 19. a 25. epoche
rozmrazí vždy jeden stupeň backbonu od posledného po prvý. S posledným
stupňom sa rozmrazí aj záverečná normalizačná vrstva (LayerNorm), ktorá
priamo napája klasifikátor. Vstupná časť (stem resp. patch embedding)
zostáva zamrznutá, rovnako ako vstupná konvolúcia ResNet50.

\begin{table}[H]
    \centering
    \caption{Postupné rozmrazovanie M3--M5 -- počet rozmrazených parametrov v jednotlivých epochách}
    \label{tab:m345-unfreeze}
    \footnotesize
    \begin{tabular}{lccc}
        \toprule
        Epocha & M3 ConvNeXt-T & M4 Swin-T & M5 CoAtNet-0 \\
        \midrule
        7.  & \texttt{stages.3} +15,47 mil. & \texttt{layers.3} +15,37 mil. & \texttt{stages.3} (transformer) +12,85 mil. \\
        13. & \texttt{stages.2} +11,11 mil. & \texttt{layers.2} +10,96 mil. & \texttt{stages.2} (transformer) +12,15 mil. \\
        19. & \texttt{stages.1} +0,99 mil.  & \texttt{layers.1} +0,97 mil.  & \texttt{stages.1} (MBConv) +1,41 mil. \\
        25. & \texttt{stages.0} +0,24 mil.  & \texttt{layers.0} +0,22 mil.  & \texttt{stages.0} (MBConv) +0,24 mil. \\
        \midrule
        Celkovo parametrov & 27,83 mil. & 27,53 mil. & 26,68 mil. \\
        \bottomrule
    \end{tabular}
\end{table}

GradCAM je pri všetkých troch modeloch aplikovaný na výstup posledného
stupňa (mapa $7\times7$ s 768 kanálmi), teda na príznaky, ktoré po
priemerovaní vstupujú do klasifikátora -- ide o rovnakú úroveň ako
\texttt{layer4} pri M1. ConvNeXt a CoAtNet majú príznaky vo formáte
$C\times H\times W$ ako bežná CNN. Swin pracuje s mapou tokenov vo formáte
$H\times W\times C$, ktorú bolo pred výpočtom GradCAM treba preusporiadať
(\texttt{reshape\_transform}). Pri Swin bola vyskúšaná aj vrstva odporúčaná
dokumentáciou knižnice pytorch-grad-cam (prvá LayerNorm posledného bloku),
tá však na tomto datasete dávala výrazné artefakty v rohoch obrazu, preto
bol zvolený výstup stupňa.

Modely M3--M5 boli trénované priamo vo Windows na rovnakej GPU. Kvôli
pamäti RAM bolo použitých 8 procesov na načítanie dát (namiesto 16). Počet
procesov ovplyvňuje iba rýchlosť, nie výsledky. Spolu s gradient
checkpointingom pri rozmrazovaní to znamená, že časy trénovania M3--M5 sú
s časmi M1 a M2 porovnateľné iba približne.

\section{Model M3 -- ConvNeXt-T}
\label{sec:m3}

\subsection{Zamrznutý backbone}

\obrazok[0.9\textwidth]{img/m3_curves.png}{M3 (zamrznutý backbone) --
priebeh trénovania najlepšieho behu (Loss a Accuracy)}{fig:m3-curves}

Už po prvej epoche má M3 validačnú accuracy $\sim$82~\% (M1 $\sim$69~\%),
príznaky ConvNeXt sú teda pre jedlá oveľa lepšie oddeliteľné lineárnym
klasifikátorom. Trénovacia accuracy na konci zostáva mierne pod validačnou
(88,6--88,9~\% oproti 89,7--90,1~\%), pretože trénovacie obrázky sú
augmentované, a validačný loss klesá až do poslednej epochy -- model sa
nepreučuje a s viac epochami by sa ešte mierne zlepšil.

\begin{table}[H]
    \centering
    \caption{M3 -- testovacie metriky (priemer $\pm$ smerodajná odchýlka z 3 behov)}
    \label{tab:m3-metrics}
    \begin{tabular}{lcc}
        \toprule
        Metrika & Zamrznutý backbone & Postupné rozmrazovanie \\
        \midrule
        Accuracy  & 89,62 $\pm$ 0,30~\% & \textbf{93,49 $\pm$ 0,36~\%} \\
        Precision & 89,82 $\pm$ 0,31~\% & \textbf{93,60 $\pm$ 0,35~\%} \\
        Recall    & 89,62 $\pm$ 0,30~\% & \textbf{93,49 $\pm$ 0,36~\%} \\
        F1        & 89,62 $\pm$ 0,30~\% & \textbf{93,50 $\pm$ 0,36~\%} \\
        Čas (3 behy) & 0:31:19 & 0:51:13 \\
        \bottomrule
    \end{tabular}
\end{table}

Jednotlivé behy so zamrznutým backbonom dosiahli 89,27~\%, 89,60~\%
a 90,00~\%. Iba s 11\,535 trénovateľnými parametrami je M3 lepší ako M1
(+3,5~p.\,b.) aj ako M2 s blokmi CBAM (+0,6~p.\,b.).

\obrazok[0.8\textwidth]{img/m3_confusion.png}{M3 (zamrznutý backbone) --
kontingenčná tabuľka najlepšieho behu (testovacie dáta)}{fig:m3-confusion}

V súčte troch behov sú najčastejšie zámeny \emph{pancakes} $\rightarrow$
\emph{waffles} (20-krát), \emph{beef\_tartare} $\rightarrow$
\emph{lasagna} (17) a \emph{chicken\_wings} $\rightarrow$ \emph{steak}
(16). Výrazne najhoršie rozpoznávaná trieda je \emph{donuts} (recall
78,3~\%), ktorá sa zamieňa hlavne s \emph{waffles} (14-krát), nasledujú
\emph{pancakes} (83,3~\%) a \emph{panna\_cotta} (86,7~\%). Najlepšie sú
\emph{greek\_salad} (97,0~\%), \emph{waffles} (95,7~\%)
a \emph{lobster\_bisque} (95,3~\%).

\subsection{Postupné rozmrazovanie}

\obrazok[0.9\textwidth]{img/m3_gradual_curves.png}{M3 (postupné
rozmrazovanie) -- priebeh trénovania najlepšieho behu (Loss
a Accuracy)}{fig:m3-gradual-curves}

\obrazok[0.9\textwidth]{img/m3_compare_gradual_unfreeze_curves.png}{M3 --
vplyv postupného rozmrazovania na priebeh validačného Loss a Accuracy
(najlepší beh)}{fig:m3-unfreeze-curves}

Po rozmrazení \texttt{stages.3} v 7. epoche stúpne validačná accuracy
z $\sim$87~\% na $\sim$92~\% a validačný loss klesne z 0,56 na
0,27--0,30. Najnižší validačný loss (0,24) dosahuje M3 v 9.--11. epoche
a najvyššiu validačnú accuracy (93,2--93,9~\%) v 9.--16. epoche. Potom sa
model začne mierne preučovať: trénovacia accuracy na konci dosiahne
99,4--99,6~\% a validačný loss vzrastie na $\sim$0,33. V porovnaní
s ostatnými modelmi je však nárast malý (kapitola~\ref{sec:m1m5}).
Postupné rozmrazovanie zlepšilo testovaciu accuracy o
\textbf{3,9~p.\,b.} na 93,49~\% (behy 93,20~\%, 94,00~\% a 93,27~\%).

\obrazok[0.8\textwidth]{img/m3_gradual_confusion.png}{M3 (postupné
rozmrazovanie) -- kontingenčná tabuľka najlepšieho behu (testovacie
dáta)}{fig:m3-gradual-confusion}

Chyby sú po rozmrazení rozložené rovnomernejšie -- žiadna zámena sa
v súčte troch behov nevyskytla viac ako 10-krát (napr. \emph{steak}
$\rightarrow$ \emph{chicken\_wings}, \emph{pancakes} $\rightarrow$
\emph{waffles}, \emph{french\_fries} $\rightarrow$ \emph{hot\_dog},
\emph{cup\_cakes} $\rightarrow$ \emph{donuts}). Trieda \emph{donuts} sa
zlepšila zo 78,3~\% na 92,3~\%. Najslabšie zostávajú \emph{pancakes}
(88,0~\%) a \emph{steak} (90,7~\%), najlepšie \emph{waffles}
a \emph{beef\_tartare} (obe 96,7~\%).

\subsection{Vizualizácia rozhodovania -- GradCAM}

\obrazok[\textwidth]{img/m3_gradcam.png}{M3 (postupné rozmrazovanie) --
GradCAM vizualizácia na správne (hore) a nesprávne (dole) klasifikovaných
testovacích vzorkách}{fig:m3-gradcam}

Pri správnych predikciách sa aktivácia sústreďuje na jedlo (mäso pri
\emph{steak}, hranolky, okraj misky s tatarákom). Pri cupcaku sa model
pozerá na papierový košíček, čo je pre túto triedu typický znak. Dve
z chýb tvoria hranolky klasifikované ako \emph{hot\_dog}: na jednom obrázku
sú hranolky podávané pri hot dogu, na druhom má jedlo tvar podlhovastej
žemle. Pri šiškách klasifikovaných ako \emph{french\_fries} obrázok
zobrazuje iba krabicu s logom, bez samotného jedla.

\section{Model M4 -- Swin-T}
\label{sec:m4}

\subsection{Zamrznutý backbone}

\obrazok[0.9\textwidth]{img/m4_curves.png}{M4 (zamrznutý backbone) --
priebeh trénovania najlepšieho behu (Loss a Accuracy)}{fig:m4-curves}

M4 má zo všetkých modelov najrýchlejší nábeh (validačná accuracy po prvej
epoche $\sim$82,5~\%) aj najnižší validačný loss na konci trénovania
($\sim$0,32). Ako pri M3, trénovacia accuracy (89,8--90,0~\%) zostáva pod
validačnou (90,3--90,5~\%) a model sa nepreučuje.

\begin{table}[H]
    \centering
    \caption{M4 -- testovacie metriky (priemer $\pm$ smerodajná odchýlka z 3 behov)}
    \label{tab:m4-metrics}
    \begin{tabular}{lcc}
        \toprule
        Metrika & Zamrznutý backbone & Postupné rozmrazovanie \\
        \midrule
        Accuracy  & 91,47 $\pm$ 0,00~\% & \textbf{94,04 $\pm$ 0,31~\%} \\
        Precision & 91,57 $\pm$ 0,02~\% & \textbf{94,17 $\pm$ 0,30~\%} \\
        Recall    & 91,47 $\pm$ 0,00~\% & \textbf{94,04 $\pm$ 0,31~\%} \\
        F1        & 91,48 $\pm$ 0,00~\% & \textbf{94,05 $\pm$ 0,31~\%} \\
        Čas (3 behy) & 0:38:30 & 1:05:41 \\
        \bottomrule
    \end{tabular}
\end{table}

So zamrznutým backbonom dosiahol M4 \textbf{najlepší výsledok spomedzi
všetkých modelov} (91,47~\%). Všetky tri behy klasifikovali správne
presne 1\,372 z 1\,500 testovacích obrázkov. Nejde o chybu -- behy majú
rôzne kontingenčné tabuľky, rôzny priebeh loss aj rôznu najlepšiu epochu
(25, 25 a 27) a mierne sa líši aj precision. Pri zamrznutom backbone sa
trénuje iba lineárna vrstva nad tými istými príznakmi, takže rôzne
inicializácie konvergujú k veľmi podobnému riešeniu a zhoda v počte
správnych odpovedí je tu pravdepodobná.

\obrazok[0.8\textwidth]{img/m4_confusion.png}{M4 (zamrznutý backbone) --
kontingenčná tabuľka najlepšieho behu (testovacie dáta)}{fig:m4-confusion}

V súčte troch behov sú najčastejšie zámeny \emph{cup\_cakes}
$\rightarrow$ \emph{donuts} (16-krát), \emph{donuts} $\rightarrow$
\emph{pancakes} (15) a \emph{french\_fries} $\rightarrow$
\emph{hot\_dog} (13). Recall všetkých tried je aspoň 87,7~\%
(\emph{pancakes}), najlepšie sú \emph{lobster\_bisque} (97,0~\%)
a \emph{greek\_salad} (95,0~\%).

\subsection{Postupné rozmrazovanie}

\obrazok[0.9\textwidth]{img/m4_gradual_curves.png}{M4 (postupné
rozmrazovanie) -- priebeh trénovania najlepšieho behu (Loss
a Accuracy)}{fig:m4-gradual-curves}

\obrazok[0.9\textwidth]{img/m4_compare_gradual_unfreeze_curves.png}{M4 --
vplyv postupného rozmrazovania na priebeh validačného Loss a Accuracy
(najlepší beh)}{fig:m4-unfreeze-curves}

Rozmrazenie \texttt{layers.3} v 7. epoche zvýši validačnú accuracy
z $\sim$88~\% na $\sim$91,5~\% a validačný loss klesne z 0,43 na
0,26--0,29. Najlepšiu validačnú accuracy (92,6--93,9~\%) dosahuje M4
v 11.--12. epoche. Po rozmrazení \texttt{layers.2} v 13. epoche sa
trénovanie stáva nestabilným: validačná accuracy kolíše v rozsahu
približne 88,5--92~\% a validačný loss rastie až na 0,42--0,49, zatiaľ
čo trénovacia accuracy dosahuje $\sim$97,5~\%. Swin sa teda pri
zvolenom learning rate $3\times10^{-4}$ preučuje výraznejšie ako ConvNeXt.
Vďaka výberu checkpointu podľa validačnej accuracy to výsledok
nepoškodilo -- testovacia accuracy stúpla o \textbf{2,6~p.\,b.} na
\textbf{94,04~\%} (behy 94,47~\%, 93,93~\% a 93,73~\%), čo je najlepší
výsledok celého projektu.

\obrazok[0.8\textwidth]{img/m4_gradual_confusion.png}{M4 (postupné
rozmrazovanie) -- kontingenčná tabuľka najlepšieho behu (testovacie
dáta)}{fig:m4-gradual-confusion}

M4 s postupným rozmrazovaním má najrovnomernejšie rozložené chyby: žiadna
zámena sa v súčte troch behov nevyskytla viac ako 8-krát
(\emph{panna\_cotta} $\rightarrow$ \emph{lobster\_bisque},
\emph{cup\_cakes} $\rightarrow$ \emph{donuts}) a recall všetkých tried je
aspoň 90,7~\% (\emph{french\_fries}). Oproti M3 sa výrazne zlepšil
\emph{steak} (z 90,7~\% na 92,7~\%) aj \emph{pancakes} (z 88,0~\% na
93,3~\%).

\subsection{Vizualizácia rozhodovania -- GradCAM}

\obrazok[\textwidth]{img/m4_gradcam.png}{M4 (postupné rozmrazovanie) --
GradCAM vizualizácia na správne (hore) a nesprávne (dole) klasifikovaných
testovacích vzorkách}{fig:m4-gradcam}

Mapy Swinu sú vďaka pozornosti v oknách ohraničené na celé objekty -- pri
hranolkách alebo steaku pokrývajú takmer celé jedlo. Pri chybách model
reaguje na zavádzajúce prvky: hranolky v žemli klasifikuje ako
\emph{hot\_dog}, jedálny lístok bez jedla ako \emph{french\_fries} (aktivácia
leží na texte) a vysokú vežu z cupcakov ako \emph{donuts}.

\section{Model M5 -- CoAtNet-0}
\label{sec:m5}

\subsection{Zamrznutý backbone}

\obrazok[0.9\textwidth]{img/m5_curves.png}{M5 (zamrznutý backbone) --
priebeh trénovania najlepšieho behu (Loss a Accuracy)}{fig:m5-curves}

M5 začína s validačnou accuracy $\sim$77~\% po prvej epoche a na konci
dosahuje $\sim$88,3~\% pri validačnom loss $\sim$0,40. Trénovacia accuracy
(89,2--89,4~\%) je len mierne nad validačnou, model sa nepreučuje.

\begin{table}[H]
    \centering
    \caption{M5 -- testovacie metriky (priemer $\pm$ smerodajná odchýlka z 3 behov)}
    \label{tab:m5-metrics}
    \begin{tabular}{lcc}
        \toprule
        Metrika & Zamrznutý backbone & Postupné rozmrazovanie \\
        \midrule
        Accuracy  & 88,24 $\pm$ 0,08~\% & \textbf{92,76 $\pm$ 0,33~\%} \\
        Precision & 88,42 $\pm$ 0,08~\% & \textbf{92,95 $\pm$ 0,28~\%} \\
        Recall    & 88,24 $\pm$ 0,08~\% & \textbf{92,76 $\pm$ 0,33~\%} \\
        F1        & 88,25 $\pm$ 0,09~\% & \textbf{92,78 $\pm$ 0,33~\%} \\
        Čas (3 behy) & 0:32:47 & 1:04:48 \\
        \bottomrule
    \end{tabular}
\end{table}

Jednotlivé behy dosiahli 88,13~\%, 88,27~\% a 88,33~\%. Hoci má CoAtNet-0
zo všetkých použitých modelov najvyššiu presnosť na ImageNet (82,4~\%), so
zamrznutým backbonom je druhý najslabší -- prekonal iba M1. Presnosť na
pôvodnej úlohe teda nepredpovedá, ako dobre sa nezmenené príznaky prenesú
na novú doménu.

\obrazok[0.8\textwidth]{img/m5_confusion.png}{M5 (zamrznutý backbone) --
kontingenčná tabuľka najlepšieho behu (testovacie dáta)}{fig:m5-confusion}

Najčastejšia zámena v súčte troch behov je \emph{lasagna} $\rightarrow$
\emph{pizza} (23-krát, najviac zo všetkých modelov), nasledujú
\emph{cup\_cakes} $\rightarrow$ \emph{donuts} (16), \emph{steak}
$\rightarrow$ \emph{beef\_tartare} (15) a \emph{cup\_cakes} $\rightarrow$
\emph{dumplings} (15). Najhoršie rozpoznávané triedy sú \emph{lasagna}
(80,7~\%), \emph{pancakes} (82,7~\%) a \emph{cup\_cakes} (84,3~\%),
najlepšie \emph{lobster\_bisque} (95,3~\%) a \emph{pizza} (93,0~\%).

\subsection{Postupné rozmrazovanie}

\obrazok[0.9\textwidth]{img/m5_gradual_curves.png}{M5 (postupné
rozmrazovanie) -- priebeh trénovania najlepšieho behu (Loss
a Accuracy)}{fig:m5-gradual-curves}

\obrazok[0.9\textwidth]{img/m5_compare_gradual_unfreeze_curves.png}{M5 --
vplyv postupného rozmrazovania na priebeh validačného Loss a Accuracy
(najlepší beh)}{fig:m5-unfreeze-curves}

Po rozmrazení transformerového stupňa \texttt{stages.3} v 7. epoche stúpne
validačná accuracy z $\sim$85~\% na $\sim$91,5~\% a validačný loss klesne
z 0,53 na $\sim$0,30. Najlepšiu validačnú accuracy (92,5--92,9~\%) dosahuje
M5 v 11.--12. epoche. Podobne ako pri Swin sa model po rozmrazení
\texttt{stages.2} preučuje -- trénovacia accuracy dosiahne $\sim$98~\%
a validačný loss vzrastie na 0,41--0,49. Postupné rozmrazovanie zlepšilo
testovaciu accuracy o \textbf{4,5~p.\,b.} na 92,76~\% (behy 92,80~\%,
93,13~\% a 92,33~\%). Po dotrénovaní teda CoAtNet dobehol väčšinu
odstupu, ale zostal za M3 aj M4.

\obrazok[0.8\textwidth]{img/m5_gradual_confusion.png}{M5 (postupné
rozmrazovanie) -- kontingenčná tabuľka najlepšieho behu (testovacie
dáta)}{fig:m5-gradual-confusion}

Najčastejšie zámeny po rozmrazení sú obojsmerné \emph{lasagna}
$\leftrightarrow$ \emph{pizza} (12 + 10-krát), \emph{chicken\_wings}
$\rightarrow$ \emph{lasagna} (12) a \emph{donuts} $\rightarrow$
\emph{panna\_cotta} (11). Najslabšie triedy sú \emph{steak} (87,3~\%)
a \emph{pancakes} (88,0~\%), najlepšie \emph{greek\_salad}
a \emph{dumplings} (obe 95,7~\%).

\subsection{Vizualizácia rozhodovania -- GradCAM}

\obrazok[\textwidth]{img/m5_gradcam.png}{M5 (postupné rozmrazovanie) --
GradCAM vizualizácia na správne (hore) a nesprávne (dole) klasifikovaných
testovacích vzorkách}{fig:m5-gradcam}

Pri správnych predikciách sa aktivácia drží na jedle (šalát, tatarák,
poleva cupcaku). Nesprávne klasifikované vzorky sú väčšinou obrázky bez
jasne viditeľného jedla: jedálny lístok (\emph{dumplings} $\rightarrow$
\emph{donuts}), obal s logom (\emph{cup\_cakes} $\rightarrow$
\emph{waffles}) alebo biely dezert prekrytý ozdobou (\emph{donuts}
$\rightarrow$ \emph{panna\_cotta}). Podobné nejasné obrázky (napr. ten istý
jedálny lístok) sa medzi chybami objavujú aj pri M3 a M4 -- ide
o vzorky, ktoré sú problematické už v samotnom datasete (Food-101 obsahuje
aj nesprávne označené alebo nejednoznačné obrázky).

\section{Porovnanie modelov M1--M5}
\label{sec:m1m5}

Tabuľka~\ref{tab:m1m5} porovnáva všetkých päť modelov v dvoch
konfiguráciách, ktoré boli natrénované pre každý model. Na overenie, či
sa modely vôbec líšia, bola pre každú konfiguráciu použitá
\textbf{jednocestná ANOVA} nad testovacou accuracy troch behov každého
modelu. Dvojice modelov boli porovnané \textbf{Welchovým t-testom}
(tabuľka~\ref{tab:m1m5-ttest}).

\begin{table}[H]
    \centering
    \caption{Porovnanie M1--M5 -- testovacia accuracy v~\% (priemer $\pm$ std z 3 behov), počet parametrov a čas trénovania}
    \label{tab:m1m5}
    \footnotesize
    \begin{tabular}{lccccccc}
        \toprule
         & & \multicolumn{3}{c}{Zamrznutý backbone} & \multicolumn{3}{c}{Postupné rozmrazovanie} \\
        \cmidrule(lr){3-5} \cmidrule(lr){6-8}
        Model & Parametre & Trénovateľné & Accuracy & Čas & Accuracy & $p$ vs M1 & Čas \\
        \midrule
        M1 ResNet50      & 23,54 mil. & 30\,735  & 86,13 $\pm$ 0,16 & 0:30 & 92,60 $\pm$ 0,19 & --    & 0:44 \\
        M2 ResNet50+CBAM & 24,24 mil. & 731\,527 & 89,07 $\pm$ 0,67 & 1:06 & 91,91 $\pm$ 0,85 & 0,369 & 1:37 \\
        M3 ConvNeXt-T    & 27,83 mil. & 11\,535  & 89,62 $\pm$ 0,30 & 0:31 & 93,49 $\pm$ 0,36 & 0,054 & 0:51 \\
        M4 Swin-T        & 27,53 mil. & 11\,535  & \textbf{91,47 $\pm$ 0,00} & 0:39 & \textbf{94,04 $\pm$ 0,31} & 0,008 & 1:06 \\
        M5 CoAtNet-0     & 26,68 mil. & 11\,535  & 88,24 $\pm$ 0,08 & 0:33 & 92,76 $\pm$ 0,33 & 0,600 & 1:05 \\
        \bottomrule
    \end{tabular}
\end{table}

\begin{table}[H]
    \centering
    \caption{Párové Welchove t-testy testovacej accuracy ($p$-hodnoty, 3 behy na model); tučne $p < 0{,}05$}
    \label{tab:m1m5-ttest}
    \footnotesize
    \begin{tabular}{lccccc@{\hspace{2em}}lccccc}
        \toprule
        \multicolumn{6}{c}{Zamrznutý backbone} & \multicolumn{6}{c}{Postupné rozmrazovanie} \\
        \cmidrule(lr){1-6} \cmidrule(lr){7-12}
           & M1 & M2 & M3 & M4 & M5 &    & M1 & M2 & M3 & M4 & M5 \\
        \midrule
        M1 & --   & \textbf{0,020} & \textbf{$<$0,001} & \textbf{$<$0,001} & \textbf{$<$0,001} & M1 & --    & 0,369 & 0,054 & \textbf{0,008} & 0,600 \\
        M2 &      & --    & 0,368 & \textbf{0,037} & 0,223 & M2 &       & --    & 0,103 & 0,057 & 0,293 \\
        M3 &      &       & --    & \textbf{0,013} & \textbf{0,017} & M3 &       &       & --    & 0,176 & 0,102 \\
        M4 &      &       &       & --    & \textbf{$<$0,001} & M4 &       &       &       & --    & \textbf{0,016} \\
        \bottomrule
    \end{tabular}
\end{table}

\noindent ANOVA potvrdila rozdiel medzi modelmi v oboch konfiguráciách --
pri zamrznutom backbone veľmi výrazne ($F = 66{,}7$, $p = 3{,}6\cdot10^{-7}$),
pri postupnom rozmrazovaní slabšie ($F = 6{,}3$, $p = 0{,}009$). Pri iba
troch behoch na model majú párové testy malú silu, preto nevýznamný
rozdiel neznamená, že modely sú rovnako dobré, iba že rozdiel sa nedal
spoľahlivo preukázať.

\obrazok[0.95\textwidth]{img/m1m5_accuracy.png}{M1--M5 -- testovacia
accuracy v oboch konfiguráciách (priemer 3 behov $\pm$ std)}{fig:m1m5-acc}

\obrazok[\textwidth]{img/m1m5_metrics.png}{M1--M5 -- Accuracy, Precision,
Recall a F1 v oboch konfiguráciách (makro priemer, 3 behy $\pm$
std)}{fig:m1m5-metrics}

\obrazoknasirku{img/m1m5_val_curves.png}{M1--M5 -- priebeh validačného
Loss a Accuracy najlepšieho behu, hore zamrznutý backbone, dole postupné
rozmrazovanie}{fig:m1m5-curves}

\obrazok[\textwidth]{img/m1m5_class_recall.png}{M1--M5 -- recall
jednotlivých tried v oboch konfiguráciách, súčet troch behov (testovacie
dáta)}{fig:m1m5-recall}

\obrazok[0.8\textwidth]{img/m1m5_acc_vs_time.png}{M1--M5 -- testovacia
accuracy vs. čas trénovania troch behov}{fig:m1m5-time}

\subsection{Diskusia}

\paragraph{Kvalita predtrénovaných príznakov (zamrznutý backbone).}
Pri zamrznutom backbone sa trénuje iba lineárny klasifikátor, výsledok
preto priamo meria, ako dobre sú predtrénované príznaky použiteľné pre
jedlá. Poradie je M4 Swin-T (91,47~\%) $>$ M3 ConvNeXt-T (89,62~\%)
$\approx$ M2 (89,07~\%) $>$ M5 CoAtNet-0 (88,24~\%) $>$ M1 ResNet50
(86,13~\%). M4 je štatisticky významne lepší ako všetky ostatné modely.
Modernejšie architektúry M3 a M4 s iba 11\,535 trénovateľnými parametrami
prekonali aj M2, ktorý cez bloky CBAM trénuje 731 tisíc parametrov. CoAtNet
je napriek najvyššej presnosti na ImageNet druhý najslabší. Je to v súlade
s poznatkom, že vyššia presnosť na pôvodnej úlohe automaticky neznamená
lepšie prenosné príznaky.

\paragraph{Postupné rozmrazovanie.} Doladenie backbonu zlepšilo všetky
modely (M1 +6,5, M2 +2,8, M3 +3,9, M4 +2,6, M5 +4,5~p.\,b.) a zároveň
zmenšilo rozdiely medzi nimi: rozpätie výsledkov kleslo z 5,3~p.\,b. na
2,1~p.\,b. Najlepšie dopadol opäť \textbf{M4 Swin-T (94,04~\%)}, nasledovaný
M3 ConvNeXt-T (93,49~\%). Rozdiel medzi nimi nie je štatisticky významný
($p = 0{,}18$). Významne lepší je M4 oproti M1 ($p = 0{,}008$) a M5
($p = 0{,}016$). CoAtNet sa po doladení dostal na úroveň M1 ($p = 0{,}60$),
jeho hybridná architektúra teda na tomto datasete neprinesla výhodu.

\paragraph{Priebeh trénovania a stabilita.} Pri zamrznutom backbone sa
žiadny model nepreučuje. Modely M3--M5 majú oveľa rýchlejší nábeh ako
ResNet50 (validačná accuracy po prvej epoche 77--83~\% oproti
$\sim$69~\%). Pri postupnom rozmrazovaní je u všetkých modelov vidieť
skok v 7. epoche po rozmrazení posledného stupňa. Potom sa modely
preučujú, každý však v inej miere (Obr.~\ref{fig:m1m5-curves}).
ConvNeXt (M3) má po celý čas najnižší validačný loss a jeho nárast je
najmenší ($\sim$0,33 na konci). Transformerové modely M4 a M5 dosahujú
najlepšiu validačnú accuracy už v 11.--12. epoche a potom ich validačný
loss narastie na 0,42--0,49. Vyplýva z toho, že learning rate
$3\times10^{-4}$, ktorý je spoločný pre všetky modely, je na doladenie
transformerov príliš vysoký. Pre praktické použitie by bolo vhodné pre M4
a M5 znížiť learning rate pre rozmrazené vrstvy (napr. na $10^{-4}$) alebo
skrátiť trénovanie. V projekte to nebolo urobené, aby zostali parametre
trénovania pre všetky modely rovnaké. Smerodajná odchýlka medzi behmi je
pri M3--M5 0,0--0,4~p.\,b., teda podobne nízka ako pri M1 a výrazne nižšia
ako pri M2.

\paragraph{Triedy.} Heatmapa recallu (Obr.~\ref{fig:m1m5-recall}) ukazuje,
že ťažké triedy sú pri všetkých modeloch podobné: \emph{pancakes}
(zámena s \emph{waffles}), \emph{steak} (zámena s \emph{beef\_tartare}
a \emph{chicken\_wings}), dvojica \emph{lasagna}--\emph{pizza}
a \emph{cup\_cakes}--\emph{donuts}. Naopak \emph{greek\_salad}
a \emph{lobster\_bisque} sú pri všetkých modeloch medzi najlepšími. Po
postupnom rozmrazení majú všetky modely recall každej triedy aspoň
87~\%, najvyrovnanejší je M4 (90,7--97,3~\%).

\paragraph{Výpočtová náročnosť.} Graf na Obr.~\ref{fig:m1m5-time} ukazuje,
že najlepší pomer presnosti a času má M3 s postupným rozmrazovaním
(93,5~\% za 51 minút). M4 je o 0,55~p.\,b. lepší, ale trénuje sa
o $\sim$28~\% dlhšie. M2 je najpomalší a so zamrznutým backbonom aj
s rozmrazovaním zaostáva za M3 aj M4. Časy M3--M5 sú s M1 a M2
porovnateľné iba približne (kapitola~\ref{sec:m345-postup}).

\section{Zhrnutie výsledkov}
\label{sec:zhrnutie}

\begin{table}[H]
    \centering
    \caption{Súhrnné porovnanie všetkých natrénovaných konfigurácií (priemer $\pm$ std z 3 behov, testovacie dáta)}
    \label{tab:summary}
    \footnotesize
    \begin{tabular}{llcccc}
        \toprule
        Model & Konfigurácia & Accuracy & Precision & Recall & Čas (3 behy) \\
        \midrule
        M1 & základná (s augm.)     & 86,13 $\pm$ 0,16 & 86,43 $\pm$ 0,15 & 86,13 $\pm$ 0,16 & 0:29:50 \\
        M1 & bez augmentácie        & 84,31 $\pm$ 0,06 & 84,67 $\pm$ 0,08 & 84,31 $\pm$ 0,06 & 0:27:14 \\
        M1 & postupné rozmraz.      & 92,60 $\pm$ 0,19 & 92,75 $\pm$ 0,15 & 92,60 $\pm$ 0,19 & 0:44:10 \\
        \midrule
        M2 & základná (s augm.)     & 89,07 $\pm$ 0,67 & 89,33 $\pm$ 0,64 & 89,07 $\pm$ 0,67 & 1:06:25 \\
        M2 & bez augmentácie        & 85,53 $\pm$ 0,52 & 86,01 $\pm$ 0,51 & 85,53 $\pm$ 0,52 & 1:04:46 \\
        M2 & postupné rozmraz.      & 91,91 $\pm$ 0,85 & 92,05 $\pm$ 0,79 & 91,91 $\pm$ 0,85 & 1:36:38 \\
        \midrule
        M3 & základná (s augm.)     & 89,62 $\pm$ 0,30 & 89,82 $\pm$ 0,31 & 89,62 $\pm$ 0,30 & 0:31:19 \\
        M3 & postupné rozmraz.      & 93,49 $\pm$ 0,36 & 93,60 $\pm$ 0,35 & 93,49 $\pm$ 0,36 & 0:51:13 \\
        \midrule
        M4 & základná (s augm.)     & 91,47 $\pm$ 0,00 & 91,57 $\pm$ 0,02 & 91,47 $\pm$ 0,00 & 0:38:30 \\
        M4 & postupné rozmraz.      & \textbf{94,04 $\pm$ 0,31} & \textbf{94,17 $\pm$ 0,30} & \textbf{94,04 $\pm$ 0,31} & 1:05:41 \\
        \midrule
        M5 & základná (s augm.)     & 88,24 $\pm$ 0,08 & 88,42 $\pm$ 0,08 & 88,24 $\pm$ 0,08 & 0:32:47 \\
        M5 & postupné rozmraz.      & 92,76 $\pm$ 0,33 & 92,95 $\pm$ 0,28 & 92,76 $\pm$ 0,33 & 1:04:48 \\
        \bottomrule
    \end{tabular}
\end{table}

Najlepší výsledok dosiahol \textbf{vizuálny transformer M4 Swin-T
s postupným rozmrazovaním (94,04~\%)}. M4 je najlepší aj so zamrznutým
backbonom (91,47~\%), jeho predtrénované príznaky sa teda na
klasifikáciu jedál prenášajú najlepšie. Tesne za ním je optimalizovaná CNN
M3 ConvNeXt-T (93,49~\%), ktorá sa od M4 štatisticky významne nelíši,
trénuje sa rýchlejšie a doladenie je pri nej najstabilnejšie. Hybridný
model M5 CoAtNet-0 ani attention vrstva CBAM v M2 po doladení neprekonali
základný ResNet50. Pri všetkých piatich modeloch platí, že postupné
rozmrazovanie je najúčinnejšie opatrenie (+2,6 až +6,5~p.\,b.) a pri M1
a M2 má väčší vplyv ako augmentácia.

\section{Implementácia}
\label{sec:implementacia}

Projekt je napísaný v Pythone (PyTorch, torchvision, timm, pytorch-grad-cam,
pandas, matplotlib, scikit-learn, SciPy). Kód v priečinku \texttt{src/} je
štruktúrovaný takto:

\begin{table}[H]
    \centering
    \caption{Štruktúra zdrojových súborov}
    \label{tab:struktura}
    \footnotesize
    \begin{tabular}{ll}
        \toprule
        Súbor & Obsah \\
        \midrule
        \texttt{config.py} & centrálna konfigurácia (cesty, split, hyperparametre) \\
        \texttt{data\_preprocessing/split\_dataset.py} & rozdelenie datasetu 80/10/10 \\
        \texttt{data\_preprocessing/dataset.py} & \texttt{Dataset}/\texttt{DataLoader} z manifestu \\
        \texttt{data\_preprocessing/transforms.py} & augmentované a neaugmentované transformácie \\
        \texttt{models/M1\_CNN\_ResNet50.py} & architektúra M1 \\
        \texttt{models/M2\_CNN\_ResNet50\_CBAM.py} & architektúra M2 (ResNet50 + CBAM) \\
        \texttt{models/timm\_backbone.py} & spoločné vytvorenie, zmrazenie a rozmrazovanie modelov z timm \\
        \texttt{models/M3\_CNN\_ConvNeXt\_Tiny.py} & M3: ConvNeXt-T, stupne rozmrazovania, vrstva pre GradCAM \\
        \texttt{models/M4\_ViT\_Swin\_Tiny.py} & M4: Swin-T, vrátane prevodu tokenov pre GradCAM \\
        \texttt{models/M5\_Hybrid\_CoAtNet\_0.py} & M5: CoAtNet-0, vlastná normalizácia vstupu \\
        \texttt{engine.py} & spoločná tréningová/evaluačná slučka, rozvrh rozmrazovania \\
        \texttt{experiment.py} & 3 behy, ukladanie metrík, grafov a checkpointov \\
        \texttt{evaluation.py} & metriky, kontingenčná tabuľka, štatistiky behov \\
        \texttt{plots.py} & grafy priebehu trénovania a kontingenčných tabuliek \\
        \texttt{gradcam.py} & GradCAM vizualizácia \\
        \texttt{M1\_ResNet50\_Train/train\_m1*.py} & M1: základná, bez augm., rozmrazovanie \\
        \texttt{M1\_ResNet50\_Train/gradcam\_m1.py} & GradCAM pre najlepší beh M1 \\
        \texttt{M2\_ResNet\_CBAM/train\_m2*.py} & M2: základná, bez augm., rozmrazovanie \\
        \texttt{M2\_ResNet\_CBAM/gradcam\_m2.py} & GradCAM pre najlepší beh M2 \\
        \texttt{M3\_ConvNeXt/}, \texttt{M4\_Swin/}, \texttt{M5\_CoAtNet/} & \texttt{train\_m*.py} (zamrznutý backbone), \\
         & \texttt{train\_m*\_gradual\_unfreeze.py}, \texttt{gradcam\_m*.py} \\
        \texttt{smoke\_test.py} & kontrola modelu pred trénovaním (pamäť GPU, gradienty, normalizácia) \\
        \texttt{summarize\_results.py} & súhrnná tabuľka konfigurácií každého modelu \\
        \texttt{compare\_ablations.py} & augmentácia a rozmrazovanie voči základnej konfigurácii \\
        \texttt{compare\_M1\_M2.py} & porovnanie M1 a M2, Welchov t-test \\
        \texttt{compare\_models.py} & porovnanie M1--M5, ANOVA, párové Welchove t-testy \\
        \bottomrule
    \end{tabular}
\end{table}

Všetky konfigurácie využívajú tú istú tréningovú a evaluačnú infraštruktúru
(\texttt{engine.py}, \texttt{experiment.py}). Líšia sa iba funkciou, ktorá
vytvorí model, a parametrami \texttt{augment\_train}
a \texttt{unfreeze\_schedule}, takže výsledky sú priamo porovnateľné.
Výsledky sú v priečinku \texttt{Results/} rozdelené podľa modelu:

\begin{itemize}
    \item \texttt{Results/M1\_CNN/} -- konfigurácie \texttt{M1\_cnn},
          \texttt{M1\_cnn\_no\_augmentation}, \texttt{M1\_cnn\_gradual\_unfreeze},
          porovnania v \texttt{comparisons/} a súhrn \texttt{results\_summary.md},
    \item \texttt{Results/M2\_CNN\_CBAM/} -- analogicky \texttt{M2\_cnn\_cbam},
          \texttt{M2\_cnn\_cbam\_no\_augmentation},
          \texttt{M2\_cnn\_cbam\_gradual\_unfreeze}, \texttt{comparisons/}
          a \texttt{results\_summary.md},
    \item \texttt{Results/M1\_vs\_M2/} -- porovnanie oboch modelov
          (tabuľka s t-testom a grafy),
    \item \texttt{Results/M3\_ConvNeXt/}, \texttt{Results/M4\_Swin/},
          \texttt{Results/M5\_CoAtNet/} -- konfigurácie \texttt{M3\_convnext},
          \texttt{M3\_convnext\_gradual\_unfreeze} (analogicky
          \texttt{M4\_swin*} a \texttt{M5\_coatnet*}), porovnanie
          v \texttt{comparisons/} a \texttt{results\_summary.md},
    \item \texttt{Results/M1\_to\_M5/} -- porovnanie všetkých modelov
          (\texttt{comparison\_table.md}, \texttt{stat\_tests.md} a grafy).
\end{itemize}

\noindent Umiestnenie dát a počet procesov na načítanie dát sa dajú
nastaviť premennými prostredia \texttt{HNS\_DATA\_DIR}
a \texttt{HNS\_NUM\_WORKERS}. Pri trénovaní vo Windows je vhodné
\texttt{HNS\_NUM\_WORKERS=8}, pretože každý proces si samostatne načítava
PyTorch a pri 16 procesoch na každý \texttt{DataLoader} dôjde operačná
pamäť.

\noindent Každá konfigurácia obsahuje pre každý beh (\texttt{run1}--\texttt{run3})
históriu trénovania (\texttt{history.csv}), kontingenčnú tabuľku (CSV aj PNG)
a graf priebehu trénovania. Na úrovni konfigurácie je súhrn troch behov
(\texttt{summary.json}), grafy najlepšieho behu, GradCAM vizualizácia
a kompletný výpis trénovania (\texttt{training\_log.txt}). Checkpointy
(\texttt{best.pt}, 90--110~MB)
sa kvôli veľkosti neukladajú do gitu. Dajú sa znova vytvoriť spustením
trénovacích skriptov.
